\documentclass[10pt,a4paper]{amsart}
\usepackage[margin=19mm]{geometry}
\usepackage{amsmath,amssymb,mathtools}
\usepackage[scr=rsfs,cal=euler]{mathalfa}
\usepackage{xfrac}
\usepackage[unicode,hidelinks,bookmarks=false]{hyperref}
\newcommand{\ie}{i.e.\ }
\newcommand{\cf}{cf.\ }
\newcommand{\resp}{respectively\ }
\newcommand{\Subsection}{\subsection}
\providecommand{\nobreakdash}{\nobreak\hbox{-}}
\setlength{\emergencystretch}{2em}
\multlinegap0pt
\usepackage{bm}
\usepackage{amssymb}
\usepackage{xcolor}
\usepackage{enumerate}
\newtheorem{lem}{Lemma}
\newtheorem{thm}{Theorem}
\DeclareMathOperator{\diver}{\mathrm{div}}
\title{Global Regular Solutions \\[0.5mm] of the Compressible Navier-Stokes Equations \\[0.5mm] with Nonlinear Density-Dependent Viscosities\\[0.5mm] and Large Initial Data of Spherical Symmetry}
\author{Gui-Qiang G. Chen, Jiawen Zhang,, Shengguo Zhu}
\date{Source: arXiv:2605.17121v1 (CC BY 4.0)}
\begin{document}
\maketitle

\noindent\emph{Source and licence.} Global Regular Solutions \\[0.5mm] of the Compressible Navier-Stokes Equations \\[0.5mm] with Nonlinear Density-Dependent Viscosities\\[0.5mm] and Large Initial Data of Spherical Symmetry. Version arXiv:2605.17121v1 (CC BY 4.0), distributed by arXiv under the Creative Commons Attribution 4.0 International licence, http://creativecommons.org/licenses/by/4.0/, as recorded in the arXiv licence metadata for this exact version and shown on the abstract page https://arxiv.org/abs/2605.17121v1. Full licence notice: \texttt{licenses/arxiv-2605.17121-CC-BY-4.0.txt}.\par\smallskip

\noindent\textit{Editorial scope.} Counted unit: the article's thm nltm (source0.tex lines 5257--5260). The article's complete original proof of it is delivered with it (source0.tex lines 5261--5545, one \texttt{proof} environment of 285 lines). No mathematics is written, repaired or re-derived: the statement and the proof are the article's own lines, in the article's own notation, and no retained line is substituted or re-flowed.  Retained uncounted as the source-local context the proof needs: lines 824--826; lines 828--833; lines 835--838; lines 4280--4288; lines 4312--4331; lines 4315--4322; lines 4411--4447; lines 4412--4415; lines 4441--4447; lines 5140--5152; lines 5164--5166; lines 5245--5254.  Nothing was omitted from any retained span, so the package carries no omission record.  The retained lines cite 1 works, whose entries are transcribed verbatim from the article's own bibliography source in the e-print and delivered with short editorial labels recorded in the item's reference mapping.  No retained line contains a figure environment, an image-inclusion command or drawing code, so no image is delivered and the package declares no asset dependency.  Editorial formatting only, in addition: an AMS \texttt{amsart} preamble that restates the article's own declarations and macros used by the retained lines, namely the environments \texttt{lem}, \texttt{thm} on separate counters, together with the standard packages those lines need.  The retained environments print in this document as Lemma 1, Theorem 1 in order of appearance, whereas the article prints the same items with the numbering of its own full document.  The complete native source of the article as distributed in the arXiv e-print for this exact version is bundled unchanged as \texttt{sources/200767/source0.tex}.
\begin{equation}\label{canshu}
\gamma\in (1,\infty),\qquad\delta\in (1-\frac{1}{n},1).    
\end{equation}

\begin{equation}\label{th78qq}
\begin{aligned}
&0<\phi_0^\frac{1}{\gamma-1}\in L^1(\mathbb{R}^n),\quad \nabla\phi_0\in H^2 (\mathbb{R}^n),\quad \boldsymbol{\psi}_0 \in  D^{1,n}(\mathbb{R}^n)\cap  D^2(\mathbb{R}^n),\\
&\nabla\phi_0^\iota\in L^{2n}(\mathbb{R}^n),\qquad \boldsymbol{u}_0\in H^3(\mathbb{R}^n),
\end{aligned}
\end{equation}

\begin{equation}\label{th78zxq}
\nabla \boldsymbol{u}_0=\phi^{-\iota}_0 \mathcal{G}_1,\qquad
\nabla (\phi^{2\iota}_0L\boldsymbol{u}_0)=\phi^{-\iota}_0\mathcal{G}_2,\qquad  L\boldsymbol{u}_0= \phi^{-2\iota}_0\boldsymbol{g}_*,
\end{equation}

\begin{equation}\label{ln}
\begin{cases}
\phi^{\epsilon,\eta}_t+\boldsymbol{w}\cdot\nabla\phi^{\epsilon,\eta}+(\gamma-1)\phi^{\epsilon,\eta} \diver\boldsymbol{w}=0,\\[3pt]
\boldsymbol{u}^{\epsilon,\eta}_t+\boldsymbol{w}\cdot \nabla \boldsymbol{w}+\nabla\phi^{\epsilon,\eta}+a \sqrt {(h^{\epsilon,\eta})^2+\epsilon^2} L\boldsymbol{u}^{\epsilon,\eta}=\boldsymbol{\psi}^{\epsilon,\eta} \cdot Q(\boldsymbol{w}),\\[3pt]
h^{\epsilon,\eta}_t+\boldsymbol{w}\cdot \nabla h^{\epsilon,\eta}+(\delta-1)g\diver\boldsymbol{w}=0,\\[3pt]
(\phi^{\epsilon,\eta},\boldsymbol{u}^{\epsilon,\eta},h^{\epsilon,\eta})|_{t=0}=(\phi^\eta_0,\boldsymbol{u}^\eta_0,h^\eta_0)=(\phi_0+\eta,\boldsymbol{u}_0,(\phi_0+\eta)^{2\iota})\qquad  \text{for $\boldsymbol{x}\in\mathbb{R}^n$},\\[3pt]
(\phi^{\epsilon,\eta},\boldsymbol{u}^{\epsilon,\eta},h^{\epsilon,\eta})\to (\eta,\boldsymbol{0},\eta^{2\iota}) \qquad \text{as $|\boldsymbol{x}|\to \infty$} \quad \text{for $t\in [0,T]$},
\end{cases}
\end{equation}

\begin{lem}\label{ls}
Let \eqref{canshu} hold, $\epsilon>0$, and $\eta>0$. 
If  the initial data  $(\phi_0, \boldsymbol{u}_0)(\boldsymbol{x})$ are spherically symmetric and satisfy \eqref{th78qq}{\rm--}\eqref{th78zxq}, then, for any  $T>0$,  there exists a unique classical solution $(\phi^{\epsilon,\eta},\boldsymbol{u}^{\epsilon,\eta},h^{\epsilon,\eta})$ in $[0,T]\times \mathbb{R}^n$ of the Cauchy problem $\eqref{ln}$ such that
\begin{align}
&\inf_{[0,T]\times \mathbb{R}^n}\phi(t,x)>0,\quad (\nabla \phi^{\epsilon,\eta},\phi^{\epsilon,\eta}_t)\in C([0,T];H^2(\mathbb{R}^n)), \notag\\
&h^{\epsilon,\eta}\in L^\infty([0,T]\times \mathbb{R}^n)\cap C([0,T]\times \mathbb{R}^n),\quad \nabla h^{\epsilon,\eta}\in C([0,T];H^2(\mathbb{R}^n)),\notag\\
&h^{\epsilon,\eta}_t\in C([0,T];H^2(\mathbb{R}^n)),\quad \boldsymbol{u}^{\epsilon,\eta}\in C([0,T];H^3(\mathbb{R}^n))\cap L^2([0,T];H^4(\mathbb{R}^n)),\label{linearregularity}\\
& \boldsymbol{u}^{\epsilon,\eta}_t\in C([0,T];H^1(\mathbb{R}^n))\cap L^2([0,T];D^2(\mathbb{R}^n)),\quad \boldsymbol{u}^{\epsilon,\eta}_{tt}\in L^2([0,T];L^2(\mathbb{R}^n)),\notag\\
&t^{\frac{1}{2}}\boldsymbol{u}^{\epsilon,\eta}\in L^\infty([0,T];D^4(\mathbb{R}^n)), \quad t^{\frac{1}{2}}\boldsymbol{u}^{\epsilon,\eta}_t\in L^\infty([0,T];D^2(\mathbb{R}^n))\cap L^2([0,T];D^3(\mathbb{R}^n)),\notag\\
& t^{\frac{1}{2}}\boldsymbol{u}^{\epsilon,\eta}_{tt}\in L^\infty([0,T];L^2(\mathbb{R}^n))\cap L^2([0,T];D^1(\mathbb{R}^n)).\notag
\end{align}
Moreover, $(\phi^{\epsilon,\eta},\boldsymbol{u}^{\epsilon,\eta},h^{\epsilon,\eta},\boldsymbol{\psi}^{\epsilon,\eta})$ is  spherically symmetric and takes the form
\begin{equation}\label{2..}
\begin{aligned}
(\phi^{\epsilon,\eta},h^{\epsilon,\eta})(t, \boldsymbol{x})&=(\phi^{\epsilon,\eta},\phi^{\epsilon,\eta})(t,|\boldsymbol{x}|),\quad 
(\boldsymbol{u}^{\epsilon,\eta},\boldsymbol{\psi}^{\epsilon,\eta})(t, \boldsymbol{x})=(u^{\epsilon,\eta}, \psi^{\epsilon,\eta})(t,|\boldsymbol{x}|) \frac{\boldsymbol{x}}{|\boldsymbol{x}|},
\end{aligned}
\end{equation} 
where  $\psi^{\epsilon,\eta}(t,r)=\frac{a\delta}{\delta-1}(h^{\epsilon,\eta})_r(t,r)$.
\end{lem}

\begin{align}
&\inf_{[0,T]\times \mathbb{R}^n}\phi(t,x)>0,\quad (\nabla \phi^{\epsilon,\eta},\phi^{\epsilon,\eta}_t)\in C([0,T];H^2(\mathbb{R}^n)), \notag\\
&h^{\epsilon,\eta}\in L^\infty([0,T]\times \mathbb{R}^n)\cap C([0,T]\times \mathbb{R}^n),\quad \nabla h^{\epsilon,\eta}\in C([0,T];H^2(\mathbb{R}^n)),\notag\\
&h^{\epsilon,\eta}_t\in C([0,T];H^2(\mathbb{R}^n)),\quad \boldsymbol{u}^{\epsilon,\eta}\in C([0,T];H^3(\mathbb{R}^n))\cap L^2([0,T];H^4(\mathbb{R}^n)),\\
& \boldsymbol{u}^{\epsilon,\eta}_t\in C([0,T];H^1(\mathbb{R}^n))\cap L^2([0,T];D^2(\mathbb{R}^n)),\quad \boldsymbol{u}^{\epsilon,\eta}_{tt}\in L^2([0,T];L^2(\mathbb{R}^n)),\notag\\
&t^{\frac{1}{2}}\boldsymbol{u}^{\epsilon,\eta}\in L^\infty([0,T];D^4(\mathbb{R}^n)), \quad t^{\frac{1}{2}}\boldsymbol{u}^{\epsilon,\eta}_t\in L^\infty([0,T];D^2(\mathbb{R}^n))\cap L^2([0,T];D^3(\mathbb{R}^n)),\notag\\
& t^{\frac{1}{2}}\boldsymbol{u}^{\epsilon,\eta}_{tt}\in L^\infty([0,T];L^2(\mathbb{R}^n))\cap L^2([0,T];D^1(\mathbb{R}^n)).\notag
\end{align}

Next, for any fixed $\eta\in(0,1]$, denote 
\begin{equation}\label{etainitial}
\mathcal{G}^\eta_1=(\phi^\eta_0)^{\iota}\nabla \boldsymbol{u}^\eta_0,\qquad \mathcal{G}^\eta_2=(\phi^\eta_0)^{\iota}\nabla((\phi^\eta_0)^{2\iota}L\boldsymbol{u}^\eta_0),\qquad 
\boldsymbol{g}^\eta_*=(\phi^\eta_0)^{2\iota}L\boldsymbol{u}^\eta_0.
\end{equation}
 Since
\begin{align}
&\begin{aligned}
\|\phi_0^\iota\nabla^2\phi_0 \|_{L^2}&\leq  C\big(\|\phi_0\|_{L^{n^*}}\|\phi_0\|^{-\iota}_{L^\infty}\|\nabla\boldsymbol{\psi}_0\|_{L^n}+\|\nabla \phi^\iota_0\|_{L^{2n}}\|\nabla \phi_0\|_{L^{\frac{2n}{n-1}}}\big),\notag\\
\|(\phi^\eta_0)^\iota\nabla^2\phi^\eta_0\|_{L^2}&= \big\|\phi_0^\iota\nabla^2\phi_0 \frac{\phi_0^{-\iota}}{(\phi_0+\eta)^{-\iota}}\big\|_{L^2}\leq  \|\phi_0^\iota\nabla^2\phi_0 \|_{L^2},\notag
\end{aligned}\\
&\hspace{13mm}\begin{aligned}
\|\mathcal{G}^\eta_1\|_{L^2}&=\big\|\frac{\phi_0^{-\iota}}{(\phi_0+\eta)^{-\iota}}\mathcal{G}_1\big\|_{L^2}\leq \|\mathcal{G}_1\|_{L^2},\\ 
\|\mathcal{G}^\eta_2\|_{L^2}
&=\big\|\frac{\phi_0^{-3\iota}}{(\phi_0+\eta)^{-3\iota}}(\mathcal{G}_2-\frac{\eta\nabla\phi_0^{2\iota}}{\phi_0+\eta}\phi_0^\iota Lu_0)\big\|_{L^2}\\
&\leq C(\|\mathcal{G}_2\|_{L^2}+\|\boldsymbol{\psi}_0\|_{L^\infty}\|\phi_0\|^{-\iota}_{L^\infty}\|\boldsymbol{g}^\eta_*\|_{L^2}),\\
\|\boldsymbol{g}_*^\eta\|_{L^2}
&=\big\|\frac{\phi_0^{-2\iota}}{(\phi_0+\eta)^{-2\iota}}\boldsymbol{g}_*\big\|_{L^2}\leq \|\boldsymbol{g}_*\|_{L^2},
\end{aligned}\label{initialapproximation}\\
&\hspace{3mm}\begin{aligned}
\|\nabla^2(\phi^\eta_0)^{2\iota}\|_{L^n}
&\leq C\big\|\nabla \big( \frac{\phi_0^{-2\iota+1}}{(\phi_0+\eta)^{-2\iota+1}}\nabla \phi_0^{2\iota}\big)\big\|_{L^n}
\leq  C\big(\|\nabla \boldsymbol{\psi}_0\|_{L^n}+\|\nabla\phi^{\iota}_0 \|^2_{L^{2n}}\big),\notag\\
\|\nabla^3(\phi^\eta_0)^{2\iota}\|_{L^2}
&\leq C\big\|\nabla^2 \big( \frac{\phi_0^{-2\iota+1}}{(\phi_0+\eta)^{-2\iota+1}}\nabla \phi_0^{2\iota}\big)\big\|_{L^2}\notag\\
&\leq  C\big(\|\nabla^2 \boldsymbol{\psi}_0 \|_{L^2}+\|\phi_0\|^{-\iota}_{L^\infty}(\delta_{3n}\|\nabla\phi^{\iota}_0 \|^3_{L^{2n}}+\delta_{2n}\|\nabla\phi^{\iota}_0 \|_{L^{\infty}}\|\nabla\phi^{\iota}_0 \|^2_{L^{2n}})\big)\\
&\quad +C\|\phi_0\|^{-\iota}_{L^\infty}\|\nabla\phi^{\iota}_0 \|_{L^{n^*}}\|\nabla \boldsymbol{\psi}_0 \|_{L^n},\notag    
\end{aligned}
\end{align}
we can check that there exists a constant $c_0\geq 1$ independent of $\eta$ such that
\begin{equation}\label{2.14}
\begin{aligned}
&\|\nabla\phi^\eta_0\|_{H^2}+\|\phi^\eta_0\|_{ L^\infty}+\delta_{3n}\|\phi^\eta_0-\eta\|_{ L^6}+\|\boldsymbol{u}^\eta_0\|_{H^3}+\|\nabla h^\eta_0\|_{L^\infty\cap D^{1,n}\cap D^2}+\|\nabla (h^\eta_0)^{\frac{1}{2}}\|_{L^{2n}}
\\
&+\|(\phi^\eta_0)^\iota\nabla^2\phi^\eta_0\|_{L^2}+\|(h^\eta_0)^{-1}\|_{L^\infty}+\|\mathcal{G}^\eta_1\|_{L^2}+\|\mathcal{G}^\eta_2\|_{L^2}+\|\boldsymbol{g}^\eta_*\|_{L^2}+2+\eta\leq c_0.
\end{aligned}
\end{equation}

\begin{equation}
\mathcal{G}^\eta_1=(\phi^\eta_0)^{\iota}\nabla \boldsymbol{u}^\eta_0,\qquad \mathcal{G}^\eta_2=(\phi^\eta_0)^{\iota}\nabla((\phi^\eta_0)^{2\iota}L\boldsymbol{u}^\eta_0),\qquad 
\boldsymbol{g}^\eta_*=(\phi^\eta_0)^{2\iota}L\boldsymbol{u}^\eta_0.
\end{equation}

\begin{equation}
\begin{aligned}
&\|\nabla\phi^\eta_0\|_{H^2}+\|\phi^\eta_0\|_{ L^\infty}+\delta_{3n}\|\phi^\eta_0-\eta\|_{ L^6}+\|\boldsymbol{u}^\eta_0\|_{H^3}+\|\nabla h^\eta_0\|_{L^\infty\cap D^{1,n}\cap D^2}+\|\nabla (h^\eta_0)^{\frac{1}{2}}\|_{L^{2n}}
\\
&+\|(\phi^\eta_0)^\iota\nabla^2\phi^\eta_0\|_{L^2}+\|(h^\eta_0)^{-1}\|_{L^\infty}+\|\mathcal{G}^\eta_1\|_{L^2}+\|\mathcal{G}^\eta_2\|_{L^2}+\|\boldsymbol{g}^\eta_*\|_{L^2}+2+\eta\leq c_0.
\end{aligned}
\end{equation}

\begin{align}
\big(\|\phi\|^2_{D^1\cap D^3}+\|\phi_t\|^2_{H^2}+
\|\phi_{tt}\|^2_{L^2}+\|\boldsymbol{\psi}_t\|^2_{D^1}\big)(t)+\int^t_0\big(\|\phi_{tt}\|^2_{H^1}+\|\boldsymbol{\psi}_{tt}\|^2_{L^2}\big)\,\mathrm{d}s\leq  c_5^2,\notag\\
\big(\|\nabla\sqrt{h}\|^2_{L^{2n}}+
\|\boldsymbol{\psi}\|^2_{D^{1,n}\cap D^2}\big)(t)\leq c^2_1,\quad \big(\|h_t\|_{L^\infty}^2+\|\boldsymbol{\psi}_t\|^2_{L^2}\big)(t)\leq  c_4^2,\notag\\
\frac{2}{3}\eta^{-2\iota}\!<h^{-1}(t,\boldsymbol{x})<\!c_1, \  
\big(\|\sqrt{h}\nabla \boldsymbol{u}\|^2_{L^2}+\|\boldsymbol{u}\|^2_{H^1}\big)(t)+\int^t_0\big(\|\nabla \boldsymbol{u}\|^2_{H^1}+\|\boldsymbol{u}_t\|^2_{L^2}\big)\,\mathrm{d}s\leq c_2^2,\notag\\
\big(\|\boldsymbol{u}\|_{D^2}^2+\|h\nabla^2\boldsymbol{u}\|^2_{L^2}+\|\boldsymbol{u}_t\|^2_{L^2}\big)(t)+\int^t_0\big(\|\boldsymbol{u}\|^2_{D^3}+\|h\nabla^2\boldsymbol{u}\|_{D^1}^2+\|\boldsymbol{u}_t\|^2_{D^1}\big)\,\mathrm{d}s\leq  c_3^2,\label{key1kk}\\
\big(\|\boldsymbol{u}_t\|^2_{D^1}+\|\sqrt{h}\nabla \boldsymbol{u}_t\|^2_{L^2}+\|\boldsymbol{u}\|^2_{D^3}+\|h\nabla^2\boldsymbol{u}\|^2_{D^1}\big)(t)+\int^t_0\|\boldsymbol{u}_t\|^2_{D^2}\,\mathrm{d}s\leq  c_4^2,\notag\\
\int^t_0\big(\|\boldsymbol{u}_{tt}\|^2_{L^2}+\|\boldsymbol{u}\|^2_{D^4}+\|h\nabla^2\boldsymbol{u}\|^2_{D^2}+\|(h\nabla^2\boldsymbol{u})_t\|^2_{L^2}\big)\,\mathrm{d}s\leq  c_4^2,\quad t\|h\nabla^2\boldsymbol{u}_t(t)\|^2_{L^2}\leq  c_5^2,\notag\\
t\big(\|\boldsymbol{u}_t\|^2_{D^2}+\|\boldsymbol{u}_{tt}\|^2_{L^2}+\|\boldsymbol{u}\|^2_{D^4}\big)(t)+
\int^t_0s\big(\|\boldsymbol{u}_{tt}\|_{D^1}^2+\|\boldsymbol{u}_t\|^2_{D^3}+\|\sqrt{h} \boldsymbol{u}_{tt}\|_{D^1}^2\big)\,\mathrm{d}s\leq  c_5^2.\notag
\end{align}

\begin{lem}\label{epsilon0}
Let \eqref{canshu} hold, $\eta>0$, and $(\phi_0,\boldsymbol{u}_0)$ satisfy \eqref{th78qq}--\eqref{th78zxq}. Then there exists $T^*>0$ independent of $\eta$ and a unique classical solution $(\phi^\eta,\boldsymbol{u}^\eta, h^\eta)$ in $[0,T^*]\times\mathbb{R}^n$ of  problem \eqref{ln} with  $\epsilon=0$ satisfying \eqref{linearregularity} with $T$ replaced by $T^*$, and $\boldsymbol{\psi}^\eta=\frac{a\delta}{\delta-1}\nabla h^\eta$. Moreover, the uniform estimates in \eqref{key1kk} that are independent of $\eta$ still hold for $(\phi^\eta,\boldsymbol{u}^\eta,h^\eta)$, and $C^{-1}\leq (gh^{-1})(t,\boldsymbol{x})\leq C$ in $[0,T^*]\times\mathbb{R}^n$.
\end{lem}

\begin{equation}\label{nl}
\begin{cases}
\phi^{\eta}_t+\boldsymbol{u}^{\eta}\cdot\nabla\phi^{\eta}+(\gamma-1)\phi^{\eta} \mathrm{div}\boldsymbol{u}^{\eta}=0,\\[2pt]
\displaystyle\boldsymbol{u}^{\eta}_t+\boldsymbol{u}^{\eta}\cdot \nabla \boldsymbol{u}^{\eta}+\nabla\phi^{\eta}+a h^{\eta} L\boldsymbol{u}^{\eta}= \boldsymbol{\psi}^{\eta} \cdot Q(\boldsymbol{u}^{\eta}),\\[3pt]
h^{\eta}_t+\boldsymbol{u}^\eta\cdot\nabla h^\eta+(\delta-1) h^\eta \mathrm{div} \boldsymbol{u}^{\eta}=0,\\[3pt]
(\phi^{\eta},\boldsymbol{u}^{\eta},h^{\eta})|_{t=0}=(\phi^\eta_0,\boldsymbol{u}^\eta_0,h_0^\eta)
=(\phi_0+\eta,\boldsymbol{u}_0,(\phi_0+\eta)^{2\iota})\quad \text{in $\mathbb{R}^n$},\\[3pt]
(\phi^{\eta},\boldsymbol{u}^{\eta},h^{\eta})\to (\eta,\boldsymbol{0},\eta^{2\iota}) \quad \text{as $|\boldsymbol{x}|\to \infty$} \qquad \text{for $t\geq 0$},
\end{cases}
\end{equation}

\begin{thm}\label{nltm}
Let \eqref{canshu} hold, $\eta>0$, and  $(\phi_0,\boldsymbol{u}_0)$ satisfy \eqref{th78qq}--\eqref{th78zxq}. Then there exist $T_*>0$ independent of $\eta$ and a unique classical solution 
$(\phi^\eta, \boldsymbol{u}^\eta, h^\eta)$ in $[0,T_*]\times\mathbb{R}^n$ of problem \eqref{nl} satisfying \eqref{linearregularity}, where $T_*$ is independent of $\eta$. Moreover, the $\eta$-uniform estimates in \eqref{key1kk} still hold for $(\phi^\eta,\boldsymbol{u}^\eta,h^\eta)$ with $T^*$ replaced by $T_*$. 
\end{thm}

\begin{proof}
We divide the proof into three steps.  

\smallskip
\textbf{1.} Since $(\phi_0,\boldsymbol{u}_0)$ satisfy \eqref{th78qq}--\eqref{th78zxq}, it follows from  \eqref{etainitial}--\eqref{initialapproximation} that  there exists a constant $c_0\geq 1$ independent of $\eta$ such that \eqref{2.14} holds. Next, denote $\mathrm{X}=(\Phi,\mathrm{H})^\top$, and  let $(\phi^0,\boldsymbol{u}^0,h^0)$ be the solution of the following Cauchy problem in $(0,\infty)\times\mathbb{R}^n$:
\begin{equation*} 
\begin{cases}
\displaystyle \mathrm{X}_t+\boldsymbol{u}_0\cdot \nabla \mathrm{X}=0, \\
\mathbf{U}_t-\mathrm{H}\Delta \mathbf{U}=\boldsymbol{0}, \\
(\Phi,\mathbf{U},\mathrm{H})|_{t=0}=(\phi^\eta_0,\boldsymbol{u}^\eta_0,h^\eta_0)\quad \text{in $\mathbb{R}^n$},\\
(\Phi,\mathbf{U},\mathrm{H})\to (\eta,\boldsymbol{0},\eta^{2\iota}) \quad \quad  \ \ \ \ \text{as $|\boldsymbol{x}|\to \infty$} \quad \text{for $t\geq 0$}.
\end{cases}
\end{equation*}

Choose $\bar{T}\in (0,T^*]$ small enough such that
\begin{align}
\begin{aligned}
0<(h^0)^{-1}\leq c_1,\ \
\sup_{t\in[0,\bar{T}]}\|\nabla h^0(t)\|^2_{L^\infty\cap D^{1,n}\cap D^2}&\leq c^2_1,\notag\\
\sup_{t\in[0,\bar{T}]}\|\boldsymbol{u}^0(t)\|^2_{H^1}
+\int^{\bar{T}}_0\big(\|\boldsymbol{u}^0\|^2_{D^2}+\|\boldsymbol{u}^0_t\|^2_{L^2}\big)\,\mathrm{d}t&\leq c_2^2,\notag
\end{aligned}\\
\begin{aligned}\label{it0}
\sup_{t\in[0,\bar{T}]}\big(\|\boldsymbol{u}^0\|_{D^2}^2+\|h^0\nabla^2\boldsymbol{u}^0\|^2_{L^2}+\|\boldsymbol{u}^0_t\|^2_{L^2}\big)(t)
+\int^{\bar{T}}_0\big(\|\boldsymbol{u}^0\|^2_{D^3}+\|\boldsymbol{u}^0_t\|^2_{D^1}\big)\,\mathrm{d}t\leq c_3^2,\\
\sup_{t\in[0,\bar{T}]}\big(\|\boldsymbol{u}^0_t\|^2_{D^1}+\|h^0_t\|^2_{D^1}+\|\boldsymbol{u}^0\|^2_{D^3}\big)(t)
+\int^{\bar{T}}_0\big(\|\boldsymbol{u}^0_t\|^2_{D^2}+\|\boldsymbol{u}^0\|^2_{D^4}+\|\boldsymbol{u}^0_{tt}\|^2_{L^2}\big)\,\mathrm{d}t
\leq c_4^2,\\[-2pt]
\sup_{t\in[0,\bar{T}]}\|(h_0)^{-1}h^0_t(t)\|^2_{L^{\infty}}
+\int^{\bar{T}}_0\big(\|(h^0\nabla^2\boldsymbol{u}^0)_t\|^2_{L^2}+\|h^0\nabla^2\boldsymbol{u}^0\|^2_{D^2}\big)\,\mathrm{d}t\leq c_4^2,\\[4pt]
\sup_{t\in[0,\bar{T}]}\big(\|\sqrt{h^0}\nabla \boldsymbol{u}^0_t\|^2_{L^2}+\|h^0\nabla^2\boldsymbol{u}^0\|^2_{D^1}\big)(t)\leq c^2_4,
\end{aligned}\\[1pt]
\begin{aligned} 
\sup_{t\in[0,\bar{T}]}\big(t\|\boldsymbol{u}^0_t(t)\|^2_{D^2}+t\|h^0\nabla^2\boldsymbol{u}^0_t(t)\|^2_{L^2}+t\|\boldsymbol{u}^0_{tt}(t)\|^2_{L^2}+t\|\boldsymbol{u}^0(t)\|^2_{D^4}\big)\leq  c_5^2,\\
\int^t_0\big(s\|\boldsymbol{u}^0_{tt}\|_{D^1}^2+s\|\boldsymbol{u}^0_t\|^2_{D^3}+s\|\sqrt{h^0} \boldsymbol{u}^0_{tt}\|_{D^1}^2+\|h^0_{tt}\|^2_{D^1}\big)\,\mathrm{d}s\leq  c_5^2.\notag
\end{aligned}
\end{align}
 


\smallskip
\textbf{2. Existence}. Let $(\boldsymbol{w},g)=(\boldsymbol{u}^0, h^0)$ in problem  \eqref{ln} with $\epsilon=0$. Then, according to Lemma \ref{epsilon0},
we  can obtain a unique  classical solution $(\phi^1,\boldsymbol{u}^1,h^1)$, which solves this problem locally in time. Inductively, we construct approximate sequences $(\phi^{k+1},\boldsymbol{u}^{k+1},h^{k+1})$ as follows:
given $(\boldsymbol{u}^k, h^k)$ for $k\geq 1$, define  $(\phi^{k+1},\boldsymbol{u}^{k+1},h^{k+1})$ by solving the  linear problem:
\begin{equation}\label{k+1}
\begin{cases}
\displaystyle
\phi^{k+1}_t+\boldsymbol{u}^k\cdot\nabla\phi^{k+1}+(\gamma-1)\phi^{k+1}\mathrm{div}\boldsymbol{u}^k=0,\\[2pt]
\displaystyle
\boldsymbol{u}_t^{k+1}+\boldsymbol{u}^k\cdot\nabla \boldsymbol{u}^k+\nabla\phi^{k+1}+ah^{k+1}L\boldsymbol{u}^{k+1}
=\boldsymbol{\psi}^{k+1}\cdot Q(\boldsymbol{u}^k),\\[2pt]
\displaystyle
h^{k+1}_t+\boldsymbol{u}^k\cdot \nabla h^{k+1}+(\delta-1)h^k\mathrm{div}\boldsymbol{u}^k=0,\\[2pt]
\displaystyle
(\phi^{k+1},\boldsymbol{u}^{k+1},h^{k+1})|_{t=0}=(\phi^\eta_0,\boldsymbol{u}^\eta_0,h^\eta_0)
 \qquad \text{in $\mathbb{R}^n$}, \\[2pt]
\displaystyle
(\phi^{k+1},\boldsymbol{u}^{k+1},h^{k+1})\to (\eta,\boldsymbol{0},\eta^{2\iota}) \quad\quad \ \qquad \text{as $|\boldsymbol{x}|\to \infty$} \quad \text{for $t\geq 0$},
\end{cases}
\end{equation}
where $\boldsymbol{\psi}^{k+1}=\frac{a\delta}{\delta-1}\nabla h^{k+1}$.
It follows from Lemma \ref{epsilon0} that the solution  $(\phi^{k+1},\boldsymbol{u}^{k+1},h^{k+1})$ $(k\in \mathbb{N})$ satisfies  the uniform estimates \eqref{key1kk}  and $C^{-1}\leq (h^k(h^{k+1})^{-1})(t,\boldsymbol{x})\leq C$ in $[0,\bar{T}]\times\mathbb{R}^n$.
To show the strong convergence of $(\phi^k,\boldsymbol{u}^k,h^k)$, we set
\begin{equation*}
\begin{aligned}
&\bar{\phi}^{k+1}=\phi^{k+1}-\phi^k,\ \ \bar{\boldsymbol{u}}^{k+1}=\boldsymbol{u}^{k+1}-\boldsymbol{u}^k, \ \  \bar{\boldsymbol{\psi}}^{k+1}=\boldsymbol{\psi}^{k+1}-\boldsymbol{\psi}^k,\ \  \bar{h}^{k+1}=h^{k+1}-h^k.
\end{aligned}
\end{equation*}
Then $(\phi^k,\boldsymbol{u}^k,h^k)$ solves the following problem due to \eqref{k+1}: 
\begin{equation}\label{k+3}
\begin{cases}
\displaystyle
\bar{\phi}^{k+1}_t+\boldsymbol{u}^k\cdot\nabla\bar{\phi}^{k+1}+\bar{\boldsymbol{u}}^k\cdot\nabla\phi^k+(\gamma-1)\big(\bar{\phi}^{k+1}\mathrm{div}\boldsymbol{u}^k+\phi^k\mathrm{div}\bar{\boldsymbol{u}}^k\big)=0,\\[4pt]
\displaystyle
\bar{\boldsymbol{u}}_t^{k+1}+\boldsymbol{u}^k\cdot\nabla \bar{\boldsymbol{u}}^k+\bar{\boldsymbol{u}}^k\cdot\nabla \boldsymbol{u}^{k-1}+\nabla\bar{\phi}^{k+1}
+ah^{k+1}L\bar{\boldsymbol{u}}^{k+1}+a\bar{h}^{k+1}L\boldsymbol{u}^{k}\\[2pt]
\displaystyle
\quad =\bar{\boldsymbol{\psi}}^{k+1}\cdot Q(\boldsymbol{u}^k)+\boldsymbol{\psi}^{k}\cdot Q(\bar{\boldsymbol{u}}^k),\\[2pt]
\displaystyle
\bar{h}^{k+1}_t+\boldsymbol{u}^k\cdot\nabla \bar{h}^{k+1}+\bar{\boldsymbol{u}}^k\cdot\nabla h^k+(\delta-1)\big(h^k\mathrm{div}{\boldsymbol{u}}^k-{h}^{k-1}\mathrm{div}\boldsymbol{u}^{k-1}\big)=0,\\[2pt]
\displaystyle
(\bar{\phi}^{k+1},\bar{\boldsymbol{u}}^{k+1},\bar{h}^{k+1})|_{t=0}=(0,\boldsymbol{0},0) \qquad\text{in $\mathbb{R}^n$},
\\[4pt]
\displaystyle
(\bar{\phi}^{k+1},\bar{\boldsymbol{u}}^{k+1},\bar{h}^{k+1})\to (0,\boldsymbol{0},0) \ \ \ \ \qquad\text{as $|\boldsymbol{x}|\to \infty$} \quad \text{for $t\geq 0$}.
\end{cases}
\end{equation}

Besides, by the same argument as in the proof of \cite[Lemma 3.11]{dxz}, we can obtain the following regularity for $(\bar{\phi}^{k+1},\bar{\boldsymbol{\psi}}^{k+1},\bar{h}^{k+1})$:
\begin{equation*} 
(\bar{\phi}^{k+1},\ \bar{\boldsymbol{\psi}}^{k+1}, \ \bar{h}^{k+1}) \in L^\infty([0,\bar{T}];H^2(\mathbb{R}^n)) \qquad \text{for $k\in \mathbb{N}$}.
\end{equation*}
Moreover, it follows from the relation $\boldsymbol{\psi}^{k+1}=\frac{a\delta}{\delta-1}\nabla h^{k+1}$ and $\eqref{k+1}_3$ that 
\begin{equation}\label{psibar1}
\begin{aligned}
&\,\bar{\boldsymbol{\psi}}^{k+1}_t+\nabla\big(\boldsymbol{u}^k\cdot \bar{\boldsymbol{\psi}}^{k+1}+\bar{\boldsymbol{u}}^k\cdot\boldsymbol{\psi}^k\big)+(\delta-1)\big(\bar{\boldsymbol{\psi}}^k\mathrm{div}\boldsymbol{u}^k+\boldsymbol{\psi}^{k-1}\mathrm{div}\bar{\boldsymbol{u}}^k\big)\\[2pt]
&+a\delta\big(h^k\nabla\diver \bar{\boldsymbol{u}}^k+\bar{h}^k\nabla\mathrm{div}\boldsymbol{u}^{k-1}\big)=\boldsymbol{0}.
\end{aligned}
\end{equation}

\textbf{2.1. Estimates for $\varphi^{k+1}\bar{h}^{k+1}$.} Denote $\varphi^{k}=(h^{k})^{-1}$. Then, according to $\eqref{k+1}_3$, 
\begin{equation*}
\varphi^{k+1}_t+\boldsymbol{u}^k\cdot \nabla \varphi^{k+1}+(1-\delta)h^k (\varphi^{k+1})^2\diver\boldsymbol{u}^k=0.
\end{equation*}
Multiplying $\eqref{k+3}_3$ by $2\overline{h}^{k+1}(\varphi^{k+1})^2$ and integrating over $\mathbb{R}^n$ yield
\begin{equation}\label{go64aah}
\frac{\rm d}{{\rm d}t}\|\varphi^{k+1}\bar{h}^{k+1}\|^2_{L^2}\\ 
\leq \frac{C}{\sigma}\|\varphi^{k+1}\bar{h}^{k+1}\|^2_{L^2} + \sigma \big(\|\bar{\boldsymbol{\psi}}^{k+1}\|^2_{L^2}
+ \|(\bar{\boldsymbol{u}}^k,\, \sqrt{h^k}\diver\bar{\boldsymbol{u}}^k,\,\varphi^k\bar{h}^k)\|_{L^2}^2\big).
\end{equation}


\textbf{2.2. Estimates for $(\bar{\phi}^{k+1},\bar{\boldsymbol{u}}^{k+1},\bar{\boldsymbol{\psi}}^{k+1})$.} First, it follows from $\eqref{k+3}_2$ that
\begin{equation}\label{hubarr}
\begin{aligned}
\|h^k\nabla^2\bar{\boldsymbol{u}}^k\|_{L^2}& \leq C\big(\|(\bar{\boldsymbol{u}}^k,\,\sqrt{h^k}\nabla \bar{\boldsymbol{u}}^k,\,\bar{\boldsymbol{u}}^k_t)\|_{L^2}+\|\bar{\boldsymbol{u}}^{k-1}\|_{L^2}\big)\\
&\quad +C\big(\|\sqrt{h^{k-1}}\nabla\bar{\boldsymbol{u}}^{k-1}\|_{L^2}+\|(\bar{\boldsymbol{\psi}}^{k},\,\varphi^{k}\bar{h}^{k},\,\nabla \bar{\phi}^{k})\|_{L^2}\big),
\end{aligned}
\end{equation}
where we have used the fact that $\nabla \varphi^k=-(\varphi^k)^2\nabla h^k$ and the following estimates:
\begin{equation}\label{2Dxinfaxian}
\begin{aligned}
\|\varphi^k\bar{h}^{k}\|_{L^3}
&\leq C\big(\|\varphi^k\bar{h}^{k}\|_{L^2}+\|\nabla(\varphi^k\bar{h}^{k})\|_{L^2}\big)\leq C\|(\bar{\boldsymbol{\psi}}^{k},\,\varphi^{k}\bar{h}^{k})\|_{L^2}\big),\\
\|\bar{h}^{k}L\boldsymbol{u}^{k-1}\|_{L^2}&=\|\varphi^k\bar{h}^{k} h^k L\boldsymbol{u}^{k-1}\|_{L^2}=\|(\varphi^k\bar{h}^{k}) (h^{k-1} L\boldsymbol{u}^{k-1})((h^{k-1})^{-1}h^k)\|_{L^2}\\
& \leq C\|\varphi^k\bar{h}^{k}\|_{L^3}\|h^{k-1} L\boldsymbol{u}^{k-1}\|_{L^6}\|(h^{k-1})^{-1}h^k\|_{L^\infty}.
\end{aligned}
\end{equation}


Next, multiplying \eqref{psibar1} by $2\bar{\boldsymbol{\psi}}^{k+1}$ and integrating over $\mathbb{R}^n$, combined with \eqref{2Dxinfaxian}, give that, for any $\sigma\in (0,1)$,
\begin{equation}\label{psibar}
\begin{aligned}
\frac{\mathrm{d}}{\mathrm{d}t}\|\bar{\boldsymbol{\psi}}^{k+1}\|^2_{L^2}
&\leq C\big(\|(\bar{\boldsymbol{u}}^{k},\sqrt{h^k}\nabla \bar{\boldsymbol{u}}^{k},h^k\nabla^2\bar{\boldsymbol{u}}^k,\bar{\boldsymbol{\psi}}^k,\bar{\boldsymbol{\psi}}^{k+1})\|_{L^2}+\|\varphi^k\bar{h}^{k}\|_{L^3}\big)\|\bar{\boldsymbol{\psi}}^{k+1}\|_{L^2}\\
&\leq  \frac{C}{\sigma}\|\bar{\boldsymbol{\psi}}^{k+1}\|^2_{L^2}
+\sigma \big(\|(\bar{\boldsymbol{u}}^k,\,\sqrt{h^k}\nabla \bar{\boldsymbol{u}}^k,\,\bar{\boldsymbol{u}}^k_t)\|^2_{L^2}+\|\bar{\boldsymbol{u}}^{k-1}\|^2_{L^2}\big)\\
&\quad +\sigma\big(\|\sqrt{h^{k-1}}\nabla\bar{\boldsymbol{u}}^{k-1}\|^2_{L^2}
+\|(\bar{\boldsymbol{\psi}}^{k},\,\varphi^{k}\bar{h}^{k},\,\nabla \bar{\phi}^{k})\|^2_{L^2}\big).
\end{aligned}
\end{equation}
  

For the estimates of $\bar{\phi}^{k+1}$, multiplying $\eqref{k+3}_1$ by $2\bar{\phi}^{k+1}$ and integrating over $\mathbb{R}^n$ give 
\begin{equation}\label{phibar1}
\frac{\mathrm{d}}{\mathrm{d}t}\|\bar{\phi}^{k+1}\|^2_{L^2}\leq C\|\bar{\phi}^{k+1}\|^2_{L^2}
+C\|(\bar{\boldsymbol{u}}^{k},\,\sqrt{h^k}\nabla\bar{\boldsymbol{u}}^{k})\|_{L^2}\|\bar{\phi}^{k+1}\|_{L^2}.
\end{equation}
Furthermore, applying $2\partial^{\boldsymbol{\varsigma}}_{\boldsymbol{x}}\bar{\phi}^{k+1}\partial^{\boldsymbol{\varsigma}}_{\boldsymbol{x}}$ $(|\boldsymbol{\varsigma}|=1)$ to $\eqref{k+3}_1$  and  integrating over $\mathbb{R}^n$, we obtain 
\begin{equation*}
\frac{\mathrm{d}}{\mathrm{d}t}\|\partial^{\boldsymbol{\varsigma}}_{\boldsymbol{x}}\bar{\phi}^{k+1}\|^2_{L^2}\leq C\|\bar{\phi}^{k+1}\|^2_{H^1}+C\|(\bar{\boldsymbol{u}}^{k},\,\sqrt{h^k}\nabla\bar{\boldsymbol{u}}^{k},\,h^k\nabla^2\bar{\boldsymbol{u}}^{k})\|_{L^2}\|\nabla\bar{\phi}^{k+1}\|_{L^2},
\end{equation*}
which, along with \eqref{hubarr} and \eqref{phibar1}, implies that, for $t\in [0,\bar{T}]$,
\begin{equation}\label{phibar2}
\frac{\mathrm{d}}{\mathrm{d}t}\|\bar{\phi}^{k+1}\|^2_{H^1}\leq \frac{C}{\sigma}\|\bar{\phi}^{k+1}\|^2_{H^1
}+\sigma\|(\bar{\boldsymbol{u}}^{k},\,\sqrt{h^k}\nabla\bar{\boldsymbol{u}}^{k},\,h^k\nabla^2\bar{\boldsymbol{u}}^{k})\|^2_{L^2}.
\end{equation}





We now estimate  $\bar{\boldsymbol{u}}^{k+1}$. Multiplying $\eqref{k+3}_2$ by $2\bar{\boldsymbol{u}}^{k+1}$ and integrating over $\mathbb{R}^n$ yield 
\begin{equation}\label{ubar}
\begin{aligned}
\frac{\mathrm{d}}{\mathrm{d}t}\|\bar{\boldsymbol{u}}^{k+1}\|^2_{L^2}\!+\!aa_1\|\sqrt{h^{k+1}}\nabla\bar{\boldsymbol{u}}^{k+1}\|^2_{L^2}&\leq \frac{C}{\sigma}\|\bar{\boldsymbol{u}}^{k+1}\|^2_{L^2}+C\big(\|\bar{\phi}^{k+1}\|^2_{H^1}+\|\bar{\boldsymbol{\psi}}^{k+1}\|^2_{L^2}\big)\\
&\quad +\sigma\big(\|(\bar{\boldsymbol{u}}^{k},\,\sqrt{h^{k}}\nabla\bar{\boldsymbol{u}}^{k})\|^2_{L^2}\!+\!\|\varphi^{k+1}\bar{h}^{k+1}\|^2_{L^2}\big),
\end{aligned}
\end{equation}
where we have used \eqref{2Dxinfaxian} and the fact that
\begin{equation*}
\|h^{k+1}L\boldsymbol{u}^k\|_{L^6}\leq C\|(h^k)^{-1}h^{k+1}\|_{L^\infty}\|h^{k}L\boldsymbol{u}^k\|_{L^6}.    
\end{equation*} 

Multiplying $\eqref{k+3}_2$ by $2\bar{\boldsymbol{u}}_t^{k+1}$ and integrating over $\mathbb{R}^n$ yield 
\begin{equation*} 
\begin{aligned}
&\,\frac{\rm d}{{\rm d}t}\big(a a_1\|\sqrt{h^{k+1}}\nabla\bar{\boldsymbol{u}}^{k+1}\|^2_{L^2}
+a(a_1+a_2)\|\sqrt{h^{k+1}}\mathrm{div}\bar{\boldsymbol{u}}^{k+1}\|^2_{L^2}\big)+2\|\bar{\boldsymbol{u}}^{k+1}_t\|^2_{L^2}\\
&\leq C\big\|(\bar{\boldsymbol{u}}^k,\sqrt{h^k}\nabla\bar{\boldsymbol{u}}^k,\sqrt{h^{k+1}}\nabla\bar{\boldsymbol{u}}^{k+1},\bar{\boldsymbol{\psi}}^{k+1},\varphi^{k+1}\bar{h}^{k+1})\big\|_{L^2}\|\bar{\boldsymbol{u}}^{k+1}_t\|_{L^2}\\
&\quad +C\|\bar{\phi}^{k+1}\|_{H^1} \|\bar{\boldsymbol{u}}^{k+1}_t\|_{L^2} +C\|\sqrt{h^{k+1}}\nabla\bar{\boldsymbol{u}}^{k+1}\|_{L^2}^2,
\end{aligned}
\end{equation*}
which, along with the Young inequality, leads to
\begin{equation}\label{hubar}
\begin{aligned}
&\,\frac{\rm d}{{\rm d}t}\big(a a_1\|\sqrt{h^{k+1}}\nabla\bar{\boldsymbol{u}}^{k+1}\|^2_{L^2}
+a(a_1+a_2)\|\sqrt{h^{k+1}}\mathrm{div}\bar{\boldsymbol{u}}^{k+1}\|^2_{L^2}\big)+\|\bar{\boldsymbol{u}}^{k+1}_t\|^2_{L^2}\\
&\leq C\big(\big\|(\bar{\boldsymbol{u}}^k,
\sqrt{h^k}\nabla\bar{\boldsymbol{u}}^k,\sqrt{h^{k+1}}\nabla\bar{\boldsymbol{u}}^{k+1},\bar{\boldsymbol{\psi}}^{k+1},\varphi^{k+1}\bar{h}^{k+1})\big\|_{L^2}^2+\|\bar{\phi}^{k+1}\|_{H^1}^2 \big).
\end{aligned}
\end{equation}

\textbf{2.3. Closing the estimates.} Combining \eqref{psibar} with \eqref{phibar2}--\eqref{hubar} implies 
\begin{equation}\label{allbar}
\begin{aligned}
&\,\frac{\mathrm{d}}{\mathrm{d}t}\varGamma^{k+1}(t,\upsilon_1) +a a_1\|\sqrt{h^{k+1}}\nabla\bar{\boldsymbol{u}}^{k+1}\|^2_{L^2}+\upsilon_1\|\bar{\boldsymbol{u}}_t^{k+1}\|^2_{L^2}\\
&\leq \frac{C}{\sigma}\varGamma^{k+1}(t,\upsilon_1) \!+\!C\sigma\|\bar{\boldsymbol{u}}^k_t\|_{L^2}^2\! +\!C(\sigma+\upsilon_1)\big(\|\bar{\boldsymbol{u}}^{k-1}\|^2_{L^2}\!+\!\|\sqrt{h^{k-1}}\nabla\bar{\boldsymbol{u}}^{k-1}\|^2_{L^2}\!+\!\varGamma^{k}(t,\upsilon_1)\big),
\end{aligned}
\end{equation}
where $\upsilon_1\in (0,1)$ is a constant to be determined later and
\begin{equation*}
\begin{aligned}
\varGamma^{k+1}(t,\upsilon_1)&=  \|\bar\phi^{k+1}\|_{H^1}^2+ \|\bar{\boldsymbol{\psi}}^{k+1}\|_{L^2}^2+ \|\bar{\boldsymbol{u}}^{k+1}\|_{L^2}^2+\|\varphi^{k+1}\bar{h}^{k+1}\|^2_{L^2}\\
&\quad+\upsilon_1 a a_1\|\sqrt{h^{k+1}}\nabla\bar{\boldsymbol{u}}^{k+1}\|^2_{L^2}
+\upsilon_1 a(a_1+a_2)\|\sqrt{h^{k+1}}\mathrm{div}\bar{\boldsymbol{u}}^{k+1}\|^2_{L^2}.
\end{aligned}
\end{equation*}

Set $\sigma=\upsilon_1^{\frac{3}{2}}$ and define
\begin{equation*}
\mathcal{E}^{k+1}(t,\upsilon_1):=\sup_{s\in [0,t]}\varGamma^{k+1}(s,\upsilon_1)+\int^t_0\big(aa_1\|\sqrt{h^{k+1}}\nabla\bar{\boldsymbol{u}}^{k+1}\|^2_{L^2}+\upsilon_1\|\bar{\boldsymbol{u}}^{k+1}_t\|^2_{L^2}\big)\,\mathrm{d}s.
\end{equation*}
Then it follows from \eqref{allbar}  and the Gr\"onwall inequality  that 
\begin{equation}\label{Gamma}
\mathcal{E}^{k+1}(t,\upsilon_1)\leq C\upsilon^{\frac{1}{2}}_1 e^{C\upsilon_1^{-\frac{3}{2}}t}(t+1)\big(\mathcal{E}^{k}(t,\upsilon_1)+\mathcal{E}^{k-1}(t,\upsilon_1)\big).
\end{equation}
Since $\bar{T}\in (0,1)$, we can first choose   $\upsilon_1=\bar{\upsilon}\in (0,1)$ such that
$C\bar{\upsilon}^{\frac{1}{2}}\leq \frac{1}{64}$, and then choose  $T_*\in (0,\bar{T}]$ such that
$(1+T_*)e^{C\bar{\upsilon}^{-\frac{3}{2}}T_*}\leq 4$. Then \eqref{Gamma} becomes
\begin{equation*}
\mathcal{E}^{k+1}(T_*,\bar\upsilon)\leq \frac{1}{16}\big(\mathcal{E}^{k}(T_*,\bar\upsilon)+\mathcal{E}^{k-1}(T_*,\bar\upsilon)\big),    
\end{equation*}
which yields $\sum_{k=1}^{\infty}\mathcal{E}^{k}(T_*,\bar\upsilon)<\infty$.



This, combined with the $k$-uniform estimates \eqref{key1kk}, gives that, for any $s'\in [1,3)$,
\begin{equation*}
\|(\bar\phi^{k},\bar {\boldsymbol{u}}^{k})\|_{H^{s'}}+\|(\bar{\boldsymbol{\psi}}^{k},\bar h^{k})\|_{L^\infty}\to 0\qquad \text{as $k\to\infty$},
\end{equation*}
which implies that there exist a subsequence (still denoted by $(\phi^k,\boldsymbol{u}^k,h^k ,\boldsymbol{\psi}^k)$) and limit functions  $(\phi^\eta,\boldsymbol{u}^\eta,h^\eta, \boldsymbol{\psi}^\eta)$ such that
\begin{equation*} 
\begin{aligned}
\phi^k\to \phi^\eta\qquad &\text{in $L^\infty([0,T_*];L^\infty(\mathbb{R}^n)\cap D^1(\mathbb{R}^n)\cap D^2(\mathbb{R}^n))$},\\
\boldsymbol{u}^k\to \boldsymbol{u}^\eta \qquad   &\text{in $L^\infty([0,T_*];H^{s'}(\mathbb{R}^n))$},\\
(\boldsymbol{\psi}^k,h^k)\to (\boldsymbol{\psi}^\eta,h^\eta) \qquad  &\text{in $L^\infty([0,T_*];L^\infty(\mathbb{R}^n))$}.
\end{aligned}
\end{equation*}

On the other hand,    the local estimates  \eqref{key1kk} independent of $k$ yield that there exists a subsequence (still denoted by) $(\phi^k,\boldsymbol{u}^k, h^k, \boldsymbol{\psi}^k)$ converging to the limit $(\phi^\eta,\boldsymbol{u}^\eta, h^\eta, \boldsymbol{\psi}^\eta)$ in the weak or weak* sense.
According to the lower semi-continuity of norms,  the corresponding estimates in \eqref{key1kk} still hold for  $(\phi^\eta,\boldsymbol{u}^\eta, h^\eta,\boldsymbol{\psi}^\eta)$ with $T^*$ replaced by $T_*$. Thus, we can check that  $(\phi^\eta,\boldsymbol{u}^\eta, h^\eta)$ is a weak solution of \eqref{nl} in the sense of distributions and satisfies the equations in \eqref{nl} pointwise.

\smallskip
\textbf{3. Uniqueness and time-continuity.}  Let $(\phi_1,  \boldsymbol{u}_1, h_1)$ and $(\phi_2, \boldsymbol{u}_2, h_2)$ be two classical solutions to the   Cauchy problem \eqref{nl} satisfying the  estimates in \eqref{key1kk}.

Set 
$\boldsymbol{\psi}_i=\frac{a\delta}{\delta-1}\nabla h_i$ ($i=1,2$), and 
\begin{equation*}
\begin{aligned}
\bar{h}=h_1-h_2,\qquad \bar{\phi}= \phi_1-\phi_2,\qquad \bar{\boldsymbol{u}}=\boldsymbol{u}_1-\boldsymbol{u}_2, \qquad  \bar{\boldsymbol{\psi}}=\boldsymbol{\psi}_1-\boldsymbol{\psi}_2.
\end{aligned}
\end{equation*}
Then it follows from the equations in  \eqref{nl} that
\begin{equation}\label{eq:1.2wcvb}
\begin{cases}
\displaystyle
\bar{\phi}_t+\boldsymbol{u}_1\cdot \nabla\bar{\phi}+\bar{\boldsymbol{u}} \cdot\nabla\phi_2+(\gamma-1)(\bar{\phi}\mathrm{div}\boldsymbol{u}_1 +\phi_2\mathrm{div}\bar{\boldsymbol{u}})=0,\\[3pt]
 \displaystyle
\bar{\boldsymbol{u}}_t+ \boldsymbol{u}_1\cdot\nabla \bar{\boldsymbol{u}}+\nabla \bar{\phi}+ a h_1L\bar{\boldsymbol{u}} \\[3pt]
\displaystyle
\quad =- \bar{\boldsymbol{u}} \cdot \nabla \boldsymbol{u}_2- a \bar h L\boldsymbol{u}_2+ \boldsymbol{\psi}_1 \cdot Q( \bar{\boldsymbol{u}})+\bar{\boldsymbol{\psi}} \cdot Q(\boldsymbol{u}_2),\\[3pt]
\bar{h}_t+\boldsymbol{u}_1\cdot \nabla\bar{h}+\bar{\boldsymbol{u}}\cdot\nabla h_2+(\delta-1)(\bar{h} \mathrm{div}\boldsymbol{u}_2 +h_1\mathrm{div}\bar{\boldsymbol{u}})=0,\\[3pt]
\displaystyle
(\bar{\phi},\bar{\boldsymbol{u}},\bar{h})|_{t=0}=(0,\boldsymbol{0},0) \qquad \text{in $\mathbb{R}^n$},\\[3pt]
\displaystyle
(\bar{\phi},\bar{\boldsymbol{u}},\bar{h})\to (0,\boldsymbol{0},0) \qquad\quad \ \text{as $|\boldsymbol{x}|\to \infty$} \quad \text{for $t\geq 0$}.
\end{cases}
\end{equation}

Set 
$\varGamma(t)=\|\bar\phi\|_{H^1}^2
+\|(h_1)^{-1}\bar{h}\|_{L^2}+\|(\bar{\boldsymbol{\psi}},\, \sqrt{aa_1h_1}\nabla\bar{\boldsymbol{u}},\,\sqrt{a(a_1+a_2)h_1}\diver\bar{\boldsymbol{u}},\bar{\boldsymbol{u}})\|^2_{L^2}$. 
In a similar way for  the derivation of \eqref{Gamma}, we obtain
\begin{equation*}
\varGamma'(t)+C^{-1}\|(\nabla \bar{\boldsymbol{u}},\,\bar{\boldsymbol{u}}_t)\|^2_{L^2}
\leq \bar{H}(t)\varGamma(t),
\end{equation*}
with a continuous  function $\bar{H}(t)\in L^1(0,T_*)$. Hence, it follows from the Gr\"onwall inequality that
$\bar{\phi}=0$ and $\bar{\boldsymbol{\psi}}=\bar{\boldsymbol{u}}=\boldsymbol{0}$.



Finally, the time-continuity follows from the same procedure as in Lemma \ref{ls}. Thus, the proof of Theorem \ref{nltm} is completed.
\end{proof}

\begin{thebibliography}{99}
\bibitem[dxz]{dxz} Q. Duan, Z. Xin, and S. Zhu, On regular solutions for three-dimensional full compressible Navier-Stokes equations with degenerate viscosities and far-field vacuum, \textit{Arch. Ration. Mech. Anal.} \textbf{247} (2023), Paper No. 3, 77pp.
\end{thebibliography}
\end{document}
